\chapter{知识表示与推理机制}
\chapterart{ch05}{推理机：从事实到结论}

\begin{introbox}
\textbf{【导读】}
本体区别于普通数据库的根本能力是\textbf{推理}：从显式知识自动导出隐含结论。
本章打开推理器的黑箱。我们先盘点推理器对外提供的五种标准服务，弄清楚「推理」到底在替工程师回答哪些问题；
然后深入 OWL 推理器的内核——Tableau 算法，亲手走完两次完整推演；
再环顾推理器市场上的其他流派：以 ELK 为代表的 consequence-based 算法、规则引擎的 RETE 网络、物化与查询时推理的架构之争。
后半章转向工程中离业务最近的三类机制：SWRL 业务规则怎么写、时态知识怎么打时间戳、不确定知识怎么标概率，
最后讨论怎样让推理器扛得住生产规模的数据。
第 2 章铺好的描述逻辑地基在本章全面派上用场；但理论覆盖范围大于当前 runtime，
第 9 章只在 manifest 明示的 profile 与 oracle 内复算，不把未支持能力包装成真实代码。
\end{introbox}

\chapterbinding{semantica.chapter\_packages.vol1.ch05}{partial}{blocked}
{SWRL、DL、时态与概率材料已作为来源资产保留；受限正向链场景
\texttt{OE-V1-CH05-SCN-FORWARD-CHAIN-001} 通过适配器绑定精确结论。}
{一般 SWRL/built-ins、完整 DL tableau、时态、概率、默认与非单调推理均不执行；
资产存在不等于 parser、profile 或 oracle 存在。}

\section{推理服务全家福：推理器能替你回答什么}

「推理」听起来是一个动作，实际上是一组边界清晰的\textbf{推理服务}（reasoning services）。
工程系统与推理器打交道，从来不是笼统地说「请推理一下」，
而是调用某个具体服务：把概念层次算全、给知识库做体检、判断某个类是不是注定为空、
弄清某台设备究竟属于什么类型、回答一条带隐含知识的查询。
每种服务都有明确定义的输入与输出，背后对应同一个逻辑问题的不同问法：
\textbf{知识库 KB 是否蕴涵某个结论}。
第 2 章讲过，OWL 推理是单调的——服务输出的每一条结论都是逻辑蕴涵，
新知识只会增加结论、不会推翻旧结论，因此这些输出可以放心地被下游系统消费。
本节逐一过这五种服务，每种都给出「输入—输出—工程用途」与一个制造业例子；
读完你会发现，第 9 章调度系统的每个功能点，都能对号入座到某个服务上。

\subsection{分类：把概念格架算出来}

\textbf{分类}（classification）的输入是全部 TBox 公理，
输出是\textbf{完备的子类关系格架}：任意两个命名类之间「是否蕴涵 \(\sqsubseteq\)」全部判定一遍，
把建模者没有显式写出、但逻辑上成立的父子关系补齐。
工程用途上，这是使用频率最高的服务——Protégé 中点击推理器后出现的「推理后类层次」
（inferred class hierarchy）就是分类结果；
应用程序通过 OWL API 的 \texttt{getSubClasses} 之类接口拿到的也是它。

看一个制造例子。建模者分别定义了两个类：
「高功率设备 \(\equiv\) 设备 \(\sqcap\) \(\exists\)功率.\(\ge\)15」、
「重载设备 \(\equiv\) 设备 \(\sqcap\) \(\exists\)功率.\(\ge\)30」，
两条定义之间没有写任何显式联系。
分类服务会自动补出「重载设备 \(\sqsubseteq\) 高功率设备」：
凡功率不低于 30 的，功率必然不低于 15，数值域的包含关系被推理器翻译成了类的包含关系。
建模者少写一条公理，知识库却没有少掉这条知识——这正是「显式知识最小化、隐含知识自动化」的本体工程理想。

\subsection{一致性检查：知识库的年检}

\textbf{一致性检查}（consistency checking）的输入是整个知识库（TBox + ABox），
输出是一个布尔值：知识库有没有模型？也就是说，全部公理与事实能否同时成立。
工程用途是\textbf{质量门禁}：本体每次修改提交，持续集成流水线都应当先跑一遍一致性检查，
不一致的版本禁止合入——因为对一个不一致的知识库，任何询问都会得到「真」
（逻辑上的爆炸原理：矛盾蕴涵一切），后续的所有推理服务全部失去意义。

制造例子：TBox 里有不相交公理「运行中 \(\sqcap\) 维护中 \(\sqsubseteq\) \(\bot\)」，
而两个业务系统各自写入了一条状态——MES 说 Mill\_012 在运行，维修系统说它在维护。
一致性检查立刻亮红灯。
更进一步，主流推理器（Pellet、HermiT）还提供\textbf{解释}（justification）服务：
找出导致矛盾的\textbf{最小公理集}，本例中就是那条不相交公理加两条状态断言，
让工程师不必在几千条公理里人肉排查。
顺带一提，HermiT 内核用的是 Tableau 的改良版本——超表（hypertableau）算法，
但对使用者而言服务接口完全一致。

\subsection{可满足性检查：揪出「永远为空」的类}

\textbf{可满足性检查}（satisfiability checking）的输入是一个概念表达式 C（连同 TBox），
输出是布尔值：是否存在某个模型，让 C 的实例集非空？
若不存在，C 就是\textbf{不可满足类}——无论将来灌入什么数据，它都注定是空集，
等价于 \(\bot\)。这几乎总是建模错误的信号，Protégé 会把这样的类标成红色。

制造例子：有人定义了「待检材料设备 \(\equiv\) 设备 \(\sqcap\) 材料」，
而 TBox 中早有「设备 \(\sqcap\) 材料 \(\sqsubseteq\) \(\bot\)」（设备与材料不相交）。
「待检材料设备」即不可满足——它要求一个东西同时是设备和材料，而公理禁止这种东西存在。
这里有一对容易混淆的概念，值得专门立一块牌子：

\begin{noteblock}
\textbf{不可满足 \(\ne\) 不一致。}
TBox 里存在不可满足类，本体仍然可以是一致的——只要没人给这个类创建实例，「空类」本身并不矛盾。
一旦某个个体被断言属于不可满足类，不一致才真正发生。
工程惯例是把「出现不可满足类」与「不一致」同样当作阻断性缺陷处理：
前者是潜伏的地雷，后者是已经爆炸的地雷。
\end{noteblock}

可满足性还是其他服务的\textbf{底层原语}：判断「C \(\sqsubseteq\) D」等价于判断
「C \(\sqcap\) \(\lnot\)D 是否不可满足」——分类服务正是靠成批的可满足性测试实现的。
5.2 节的 Tableau 算法，解决的核心问题就是可满足性。

\subsection{实例检查与实现：个体究竟是什么}

\textbf{实例检查}（instance checking）的输入是一个个体 a 与一个概念 C，
输出是布尔值：KB 是否蕴涵 C(a)？
把实例检查对所有个体、所有命名类成批跑一遍，求出每个个体的\textbf{最具体类型集合}，
就是\textbf{实现}（realization）服务。
工程用途：设备档案自动归类、按类型聚合统计、给前端页面挂正确的展示模板。

制造例子：ABox 只写了「设备(Lathe\_003)、功率(Lathe\_003, 15.5)」。
实现服务跑完后，Lathe\_003 的类型集合从 \{设备\} 扩充为 \{设备, 高功率设备\}——
因为 15.5 \(\ge\) 15 触发了高功率设备的等价类定义。
没有任何人写过「Lathe\_003 是高功率设备」这条数据，它是被推出来的。
这个例子将在 5.2 节被完整解剖：推理器内部到底走了哪几步才得出这个结论。

\subsection{查询回答：让隐含知识可检索}

\textbf{查询回答}（query answering）的输入是一条合取查询
（conjunctive query，工程上通常以 SPARQL 表达），
输出是所有满足查询的变量绑定——关键在于，匹配是在\textbf{蕴涵闭包}上进行的，
显式三元组与推理出的隐含三元组同等参与匹配。

制造例子：调度面板要查「位于 WorkCell\_A、状态为 Idle 的高功率设备」。
数据库里根本没有「高功率设备」这个标签字段，
但只要推理服务补全了「Lathe\_003 属于高功率设备」，这条查询就能命中它。
没有推理的 SPARQL 只能查到「写进去的」，带推理的查询能查到「成立的」——
这是本体存储与普通图数据库拉开差距的地方，第 7 章讨论存储选型时还会回到这一点。

\subsection{前向与后向：推理发生在哪个时刻}

五种服务说的是「推理器回答什么」，还有一个正交的工程维度是「推理在什么时刻发生」，
也就是\textbf{推理策略}的选择。

\textbf{前向推理}（forward chaining）由数据驱动：事实一进来，规则就开火，
新结论立即产出并可能触发更多规则，直到再也没有新结论为止（到达不动点）。
它把推理成本付在\textbf{写入时}，查询时几乎零开销，
适合「数据到达就要反应」的场景：状态告警、冲突检测、流式监控。
5.3 节的 RETE 网络与物化架构，都是前向推理的工业形态。

\textbf{后向推理}（backward chaining）由目标驱动：先有一个待证明的问题，
推理器从目标出发反向寻找「哪些规则能推出它、这些规则的前提又靠什么成立」，
一路回溯到已知事实。它把成本付在\textbf{查询时}，不为没人问的结论买单，
适合「数据巨量、问题稀少」的场景：故障诊断、根因追溯、按需校验。
Prolog 的求解过程、OWL QL 的查询重写都是后向风格。

两种策略不是互斥的信仰，而是成本分摊方式的选择：
更新频繁、查询多样选后向；查询密集、低延迟要求高选前向；
生产系统常见混合式——核心层次前向物化，长尾问题后向求解。
仓库 \texttt{ch05-reasoning/README.md} 里有一张推理类型对照表，可以与本节互相印证；
5.3.3 与 5.7 节会把这个权衡展开成完整的架构决策。

把两种策略放进同一个业务故事里更直观。
车间大屏的「冲突告警」走前向：状态三元组一落库，规则立即开火，
告警在数据到达后的毫秒级亮起——没人查询，结论也已经在那里。
维修工的「为什么这台设备被判不可用」走后向：从结论出发反查证据链——
不可用是哪条规则推的、规则前提里的「维护中」状态又是谁写入的，
一路追到原始工单；这种问题一天只问几次，预先全量推导毫无必要。

最后把五种服务串成一条\textbf{标准作业线}，这也是推理器实际的调用顺序：
先跑\textbf{一致性检查}，不一致立即止损（其余结果全不可信）；
通过后跑\textbf{分类}，把类格架算全；可满足性检测顺势完成，红色类回炉返修；
然后\textbf{实现}给全体个体归类；最后\textbf{查询回答}面向应用持续服务。
本体的持续集成流水线就是这条作业线的脚本化：
提交触发前三步，任何一步亮红灯就阻断合入——
推理服务在工程里首先不是「智能功能」，而是\textbf{质量基础设施}。

\section{Tableau 算法精讲：用「试着造世界」回答问题}

上一节把推理服务一字排开，本节钻进引擎舱，看 OWL 推理器内核的
\textbf{Tableau 算法}如何实现这些服务。
原理可以一句话说清：\textbf{尝试构造一个满足全部公理的模型；
构造成功，说明概念可满足；所有尝试都撞上矛盾，说明不可满足}。
Pellet、HermiT（及其超表变体）、FaCT++ 的内核都是这套思路。

\subsection{核心思想：归谬式的「造世界」}

Tableau 回答的基础问题是 5.1.3 的可满足性，而其他服务都能归约到它。
最典型的是蕴涵判定走\textbf{归谬法}（refutation）：
要判断「KB 是否蕴涵 Lathe\_003 需要冷却」，
就把它的反面「Lathe\_003 不需要冷却」加进知识库，看会不会矛盾——
反面无法成立，原命题即被蕴涵。

「造世界」造的是一张\textbf{完成图}（completion graph）：
节点代表个体——既包括 ABox 里的真实个体，也包括推理过程中虚构出来的匿名个体；
每个节点挂一个\textbf{标签}，记录「这个个体必须属于哪些概念」；
节点之间的边代表属性关系。
算法从初始事实出发，反复套用一小组\textbf{展开规则}，
把标签里的复合概念逐步拆解、传播、落实；
拆到某个节点的标签里同时出现「A」与「\(\lnot\)A」（或出现 \(\bot\)），这条路就\textbf{冲突}（clash）了；
拆到再没有规则可用而又没有冲突，模型就造成了。

\subsection{四条展开规则}

构成 Tableau 主体的展开规则与第 2 章的四个核心构造器一一对应，先看总表，再逐条讲直觉。

\begin{center}
\begin{tabularx}{\linewidth}{@{}>{\centering\arraybackslash}p{0.10\linewidth}p{0.34\linewidth}X@{}}
\toprule
\textbf{规则} & \textbf{触发条件} & \textbf{执行动作} \\
\midrule
\(\sqcap\) 规则 & 节点 x 标签含 C \(\sqcap\) D & 把 C 与 D 都加入 x 的标签 \\
\(\sqcup\) 规则 & 节点 x 标签含 C \(\sqcup\) D & 分支：一条支线加 C，另一条加 D，分别继续 \\
\(\exists\) 规则 & x 标签含 \(\exists\)R.C，且 x 尚无满足 C 的 R 后继 & 新建匿名节点 y，添加边 R(x, y) 与标签 C(y) \\
\(\forall\) 规则 & x 标签含 \(\forall\)R.C，且存在边 R(x, y) 而 y 标签缺 C & 把 C 加入 y 的标签 \\
\bottomrule
\end{tabularx}
\end{center}

\(\sqcap\) \textbf{规则是拆包裹}：「C 且 D」是捆在一起的承诺，拆开来两件都得兑现。
它最便宜，不分支、不造点，只做标签登记。

\(\sqcup\) \textbf{规则是岔路口}：「C 或 D」至少居其一，但算法不知道该选哪个，
只能先押一条路走下去，留好存档；走不通（冲突）再回档换另一条。
这是 Tableau 中\textbf{不确定性的唯一来源}，也是指数复杂度的来源——
每个析取都让待探索的世界数量翻倍，5.2.5 再谈推理器怎么缓解。

\(\exists\) \textbf{规则是造证人}：标签承诺「存在一个 R 关系的 C」，
而图里找不到这样的对象，算法就亲手造一个匿名节点充当\textbf{证人}（witness）。
注意触发条件里的「尚无满足 C 的 R 后继」——已有证人就不再重复制造，这个细节是终止性的伏笔之一。

\(\forall\) \textbf{规则是巡逻队}：「所有 R 邻居都得是 C」不会主动创造任何东西，
但只要 x 长出一条 R 边，巡逻队就把 C 盖到对端节点的标签上。
它的成本在于\textbf{传播}：每条新边都要补一轮检查，
\(\forall\) 与 \(\exists\) 交替嵌套时标签会沿着链条层层扩散——这是 5.7 节「\(\forall\) 比 \(\exists\) 贵」的机理。

规则表之外还有两件配套设施值得知道。
其一是\textbf{冲突判定}的完整清单：除了标签里同时出现 A 与 \(\lnot\)A、出现 \(\bot\) 之外，
\textbf{数据类型层}也能产生冲突——若某个体的功率被同时约束到「\(\ge\)30」与「\(\le\)20」，
两个数值区间交集为空，与逻辑矛盾同等致命；
判定数值区间、字符串模式这类约束是否相容的子模块叫数据类型推理器（datatype reasoner），
它在下一小节的推演里会率先出场。
其二是 \textbf{TBox 公理怎么参与}：实现中并不把所有公理粗暴地灌进每个节点标签，
而是\textbf{惰性展开}（lazy unfolding）——节点标签里出现某个命名类时，
才把该类的定义按需展开进来，避免无关公理白白膨胀标签。
这两件设施不改变四条规则的逻辑，却决定了真实推理器的速度档位。

\subsection{完整推演之一：Lathe\_003 需要冷却吗}

现在用仓库里的完整例子走一遍流程。问题：已知下面的 TBox 与 ABox，
推理器能否得出「Lathe\_003 需要冷却系统」？

\begin{codebox}
\codeline{\textcolor{cmtgray}{\#\ ==============================}}
\codeline{\textcolor{cmtgray}{\#\ 2.\ Tableau推理演示}}
\codeline{\textcolor{cmtgray}{\#\ ==============================}}
\codeblank{}
\codeline{\textcolor{cmtgray}{\#\ 问题：Lathe\_003\ 是否需要冷却系统？}}
\codeblank{}
\codeline{\textcolor{cmtgray}{\#\ TBox（类定义）：}}
\codeline{高功率设备\ \(\equiv\)\ 设备\ \(\land\)\ \(\exists\)功率.浮点数[\(\ge\)15]}
\codeline{高功率设备\ \(\sqsubseteq\)\ \(\exists\)需要冷却.冷却系统}
\codeline{设备\ \(\land\)\ 材料\ \(\sqsubseteq\)\ \(\bot\)\ \ \ （设备和材料不相交）}
\codeblank{}
\codeline{\textcolor{cmtgray}{\#\ ABox（事实）：}}
\codeline{设备(Lathe\_003)}
\codeline{功率(Lathe\_003,\ 15.5)}
\codeblank{}
\codeline{\textcolor{cmtgray}{\#\ 推理步骤：}}
\codeline{\textcolor{cmtgray}{\#\ 步骤1：匹配等价类规则}}
\codeline{\textcolor{cmtgray}{\#\ \ \ 功率(Lathe\_003,\ 15.5)，且\ 15.5\ \(\ge\)\ 15，满足"\(\exists\)功率.浮点数[\(\ge\)15]"}}
\codeline{\textcolor{cmtgray}{\#\ \ \ 结合"高功率设备\ \(\equiv\)\ 设备\ \(\land\)\ \(\exists\)功率.浮点数[\(\ge\)15]"}}
\codeline{\textcolor{cmtgray}{\#\ \ \ \(\rightarrow\)\ 推出：Lathe\_003\ 属于高功率设备}}
\codeblank{}
\codeline{\textcolor{cmtgray}{\#\ 步骤2：应用子类公理}}
\codeline{\textcolor{cmtgray}{\#\ \ \ 已知"高功率设备\ \(\sqsubseteq\)\ \(\exists\)需要冷却.冷却系统"}}
\codeline{\textcolor{cmtgray}{\#\ \ \ 且\ Lathe\_003\ 属于高功率设备}}
\codeline{\textcolor{cmtgray}{\#\ \ \ \(\rightarrow\)\ 推出：Lathe\_003\ 属于"\(\exists\)需要冷却.冷却系统"}}
\codeblank{}
\codeline{\textcolor{cmtgray}{\#\ 步骤3：应用存在规则}}
\codeline{\textcolor{cmtgray}{\#\ \ \ Lathe\_003\ 属于"\(\exists\)需要冷却.冷却系统"，模型中没有对应的冷却系统个体}}
\codeline{\textcolor{cmtgray}{\#\ \ \ \(\rightarrow\)\ 创建新个体\ cooling\_1}}
\codeline{\textcolor{cmtgray}{\#\ \ \ \(\rightarrow\)\ 添加：需要冷却(Lathe\_003,\ cooling\_1)\ 和\ 冷却系统(cooling\_1)}}
\codeblank{}
\codeline{\textcolor{cmtgray}{\#\ 结论：Lathe\_003\ 需要冷却系统\ \(\checkmark\)}}
\codeblank{}
\codeline{\textcolor{cmtgray}{\#\ Tableau算法终止条件：}}
\codeline{\textcolor{cmtgray}{\#\ 产生冲突\ \(\rightarrow\)\ 概念C不可满足}}
\codeline{\textcolor{cmtgray}{\#\ 没有规则可应用\ \(\rightarrow\)\ 概念C可满足（找到了合法模型）}}
\codeline{\textcolor{cmtgray}{\#\ 所有分支都有冲突\ \(\rightarrow\)\ 概念C不可满足}}
\codeline{\textcolor{cmtgray}{\#\ 至少有一个分支无冲突\ \(\rightarrow\)\ 概念C可满足}}
\codeblank{}
\end{codebox}
\color{black}

\assetsource{ch05}{description-logic-examples}

把片段中的三步翻译成完成图上的动作，逐步解说如下。

\textbf{步骤 1（数值触发等价类定义）}：初始完成图只有一个节点 Lathe\_003，
标签为 \{设备\}，外加一条数据属性「功率 = 15.5」。
等价类公理「高功率设备 \(\equiv\) 设备 \(\sqcap\) \(\exists\)功率.\(\ge\)15」是双向的：
右边成立就必须把左边补上。
右边是一个 \(\sqcap\)：第一个合取项「设备」已在标签里；
第二个合取项是对数据值的存在限制，推理器的\textbf{数据类型推理}模块负责判定
15.5 落在「\(\ge\)15」的取值区间内——成立。
于是「高功率设备」被加入 Lathe\_003 的标签。
注意方向：这里用的是定义的「充分性」一侧，普通子类公理 \(\sqsubseteq\) 没有这一侧，
这正是 5.7 节「等价类比子类贵」的原因——双向都要随时盯着。

\textbf{步骤 2（子类公理传播）}：标签里有了「高功率设备」，
公理「高功率设备 \(\sqsubseteq\) \(\exists\)需要冷却.冷却系统」立即生效：
凡属于左边的，必须背上右边的承诺。
「\(\exists\)需要冷却.冷却系统」进入 Lathe\_003 的标签——此刻它还只是一张欠条。

\textbf{步骤 3（\(\exists\) 规则兑现欠条）}：算法检查 Lathe\_003 的邻居，
没有任何「需要冷却」边指向一个冷却系统——欠条没人兑付，\(\exists\) 规则触发：
新建匿名节点 cooling\_1，添加边「需要冷却(Lathe\_003, cooling\_1)」与标签「冷却系统(cooling\_1)」。

此后再检查：\(\sqcap\) 拆完了，没有 \(\sqcup\)，\(\exists\) 的证人有了，\(\forall\) 无从触发，
图上也没有任何冲突——\textbf{无规则可用且无冲突，模型构造成功}。
这个世界里 Lathe\_003 确实连着一台冷却系统，而且任何满足全部公理的世界都必然如此。

最后补上归谬视角，验证「蕴涵」确实成立：把反面断言「Lathe\_003 不需要任何冷却系统」加入知识库，
它等价于给 Lathe\_003 贴上标签「\(\forall\)需要冷却.\(\lnot\)冷却系统」。
推演照旧走到步骤 3 造出 cooling\_1 后，\(\forall\) 规则的巡逻队出动：
沿着「需要冷却」边把「\(\lnot\)冷却系统」盖到 cooling\_1 头上——
它的标签里于是同时出现「冷却系统」与「\(\lnot\)冷却系统」，冲突！
反面无路可走，原结论「Lathe\_003 需要冷却」被正式蕴涵。
两个视角殊途同归，建议读者对照片段把两遍都手算一次，体会「构造模型」与「驳倒反面」是同一枚硬币的两面。

还有一个容易被忽略的落库问题：推演结束后，哪些结论可以写回知识库？
「Lathe\_003 属于高功率设备」「Lathe\_003 属于 \(\exists\)需要冷却.冷却系统」
都是被蕴涵的类型断言，可以放心物化（5.7.4）。
但 cooling\_1 \textbf{不行}——它是推理器为证明可满足性搭的脚手架，
只承诺「存在这么一台冷却系统」，并没有说它是台账里的哪一台；
把它当作命名个体落库，等于凭空给车间登记了一台不存在的设备。
脚手架与资产的这条分界线，在 5.4.3 讨论规则可见性时还会再次出现。

\subsection{完整推演之二：一台「不在任何工作单元」的设备}

第二个例子展示 Tableau 的另一面：报告\textbf{不一致}。
包内教学资产用一阶逻辑记法给出了一对互相矛盾的公理——
为了可读性写成 FOL，但矛盾结构与 DL 表达完全相同。

\begin{codebox}
\codeline{\textcolor{cmtgray}{\#\ ==============================}}
\codeline{\textcolor{cmtgray}{\#\ 3.\ 一致性检查示例}}
\codeline{\textcolor{cmtgray}{\#\ ==============================}}
\codeblank{}
\codeline{\textcolor{cmtgray}{\#\ 一致性公理（相互矛盾的例子）：}}
\codeline{\textcolor{cmtgray}{\#\ 公理1：所有设备都必须位于某个工作单元}}
\codeline{\(\forall\)x:\ Equipment(x)\ \(\rightarrow\)\ locatedIn(x,\ WorkCell)}
\codeblank{}
\codeline{\textcolor{cmtgray}{\#\ 公理2：存在不位于任何工作单元的设备}}
\codeline{\(\exists\)x:\ Equipment(x)\ \(\land\)\ \(\lnot\)(\(\exists\)w:\ locatedIn(x,\ w))}
\codeline{\textcolor{cmtgray}{\#\ \(\uparrow\)\ 这两条公理相互矛盾，违反一致性要求}}
\codeline{\textcolor{cmtgray}{\#\ Tableau算法会检测到所有分支都产生冲突，报告不一致}}
\codeblank{}
\end{codebox}
\color{black}

\assetsource{ch05}{description-logic-examples}

推演过程如下。公理 2 是一个存在断言：天底下有一台设备不位于任何工作单元。
Tableau 把这个「存在」落实为匿名个体 e，标签登记两件事：
「Equipment(e)」与「e 没有任何 locatedIn 关系对象」——
后者等价于给 e 贴上「\(\forall\)locatedIn.\(\bot\)」式的封条：谁敢当 e 的位置，谁就是矛盾。
接着公理 1 上场：所有设备都必须位于某个工作单元，
于是 e 的标签里被塞进一张「\(\exists\)locatedIn」的欠条。
\(\exists\) 规则兑现欠条，造出位置证人 w 并连边 locatedIn(e, w)；
公理 2 的封条立即沿这条边传播 \(\bot\) 给 w——冲突。
这里没有 \(\sqcup\)，没有别的分支可试，\textbf{所有路径都撞墙，算法报告不一致}。

工程上的要点是：真实项目里的矛盾很少像教科书这样面对面站好，
它们通常隐藏在继承链与属性约束的交叠处——
例如某个子类从两个父类分别继承了「必须有位置」与「禁止有位置」。
靠人眼看不出来，靠一致性检查一秒现形，再用解释服务把肇事公理揪出来。
这也是为什么成熟团队都把本体纳入版本管理与持续集成：
\textbf{一致性检查就是本体的单元测试}。

\subsection{分支、回溯与阻塞：算法为什么会终止}

前两个例子都没有用到 \(\sqcup\) 规则，但真实本体充满析取
（「设备状态是运行、空闲或维护之一」），分支不可避免。
\textbf{分支}（branching）的代价在于：每个 \(\sqcup\) 都把当前世界复制成两份，
嵌套的析取让世界数量指数膨胀。
\textbf{回溯}（backtracking）是配套机制：某条支线冲突后，
撤销该支线的全部增改，退回最近的岔路口换下一条。
现代推理器在此之上做了\textbf{依赖导向回溯}（backjumping）：
给每个结论记账——它依赖了哪些分支选择；
冲突发生时直接跳回\textbf{真正肇事}的那个岔路口，
中间无关的选择点整段跳过，省掉成片的无效探索。
这套机账制度也是推理器性能波动的来源之一：同一个本体，公理顺序不同、版本不同，耗时可能差出数量级。

还剩最后一个疑问：\(\exists\) 规则会造新节点，新节点又可能触发新的 \(\exists\)——
算法凭什么不会永远造下去？考虑一对循环公理：
「设备 \(\sqsubseteq\) \(\exists\)属于.产线」、「产线 \(\sqsubseteq\) \(\exists\)包含.设备」。
从一台设备出发：造产线，产线要求包含设备，再造设备，新设备又要求属于产线……
朴素展开确实永不停机。
解药是\textbf{阻塞}（blocking）：
\textbf{若新节点 y 的标签是其某个祖先 x 标签的子集，就冻结 y，不再对它套用造点规则}。
直觉很简单：y 面临的「待办事项」x 全都有，x 那里怎么解决的，y 照抄即可——
与其继续展开，不如把 y 的后续需求「绕回」x 已经造好的结构，
无限延伸的链条被折叠成有限的环。
上面的例子里，第二台匿名设备的标签与第一台设备完全相同，立即被阻塞，算法停机并报告可满足。
终止性的完整论证依赖两个事实：标签只能从知识库的有限子概念集合里取值，
而阻塞限制了完成图的深度——有限标签 \(\times\) 有限深度，步数必然有限。
直觉到此足够；需要严格证明的读者可以查描述逻辑手册，工程上记住一句话：
\textbf{阻塞是 Tableau 在「表达力」与「保证停机」之间的安全阀}。

\subsection{工程中的描述逻辑约束}

理解了机理，回头看工程上最常写的四类 DL 约束——它们将成为推理器日常咀嚼的口粮。

\begin{codebox}
\codeline{\textcolor{cmtgray}{\#\ ==============================}}
\codeline{\textcolor{cmtgray}{\#\ 4.\ 工程中的描述逻辑约束}}
\codeline{\textcolor{cmtgray}{\#\ ==============================}}
\codeblank{}
\codeline{\textcolor{cmtgray}{\#\ 数控车床必须有数控系统}}
\codeline{CNCLathe\ \(\sqsubseteq\)\ \(\exists\)hasControl.CNCControl}
\codeblank{}
\codeline{\textcolor{cmtgray}{\#\ 设备利用率在0-1之间}}
\codeline{Equipment\ \(\sqsubseteq\)\ \(\forall\)hasUtilization.(\(\ge\)0\ \(\sqcap\)\ \(\le\)1.0)}
\codeblank{}
\codeline{\textcolor{cmtgray}{\#\ 正在运行的设备不能同时维护（不相交）}}
\codeline{Running\ \(\sqcap\)\ Maintenance\ \(\sqsubseteq\)\ \(\bot\)}
\codeblank{}
\codeline{\textcolor{cmtgray}{\#\ 加工设备必须能加工至少一种材料}}
\codeline{ProcessingEquipment\ \(\sqsubseteq\)\ \(\exists\)canProcess.Material}
\codeblank{}
\end{codebox}

\assetsource{ch05}{description-logic-examples}

逐条点评。
第一条「CNCLathe \(\sqsubseteq\) \(\exists\)hasControl.CNCControl」是\textbf{存在型约束}，最便宜也最常用；
但要当心开放世界语义：某台数控车床的 ABox 里没写数控系统，推理器并\textbf{不会报错}，
而是默默虚构一个匿名证人替它圆场。
想把「漏配数据」查出来，要靠封闭世界检查（SPARQL 的 \texttt{FILTER NOT EXISTS}，见第 2 章 OWA 的讨论）——
\textbf{OWL 公理是逻辑承诺，不是数据校验器}，这是新手最常踩的坑。
第二条利用率约束用 \(\forall\) 把取值锁进 0 到 1.0 的区间，属于\textbf{传播型约束}，
每条 hasUtilization 边都要被巡逻一次。
第三条「Running \(\sqcap\) Maintenance \(\sqsubseteq\) \(\bot\)」是\textbf{不相交约束}，
它是一致性检查的主要火力来源，也是 5.6 节「确定性硬骨架」的标准样本。
第四条加工能力约束与第一条同型。
四条约束的「推理价格」并不相同——\(\exists\) 便宜、\(\forall\) 贵、不相交引发冲突检测，
这份价目表在 5.7 节展开成完整的性能工程经验。

\section{其他推理风格：没有银弹的推理器市场}

Tableau 是表达力最全的通用解法，却不是唯一解法。
推理器市场长期存在几个并行流派，各自押注不同的「表达力—速度」平衡点。
了解它们，才能在选型时读懂产品说明书上那些缩写。

\subsection{Consequence-based 推理：ELK 为什么快}

Tableau 的思路是「试图造出模型」，
\textbf{consequence-based}（基于推论）算法反其道而行：
不造模型，直接按一组完备的推导规则\textbf{正向饱和}（saturation）——
从已有公理机械地推出所有能推出的子类关系，推到不动点为止，
分类结果一次性全部产出。
没有分支、没有回溯、没有匿名节点的反复构造，
推导过程还能按公理分片并行。

这套打法的代表是 \textbf{ELK} 推理器，它服务于 OWL 2 的 \textbf{EL Profile}：
一个刻意削减过的语言子集——保留 \(\sqcap\)、\(\exists\)、类层次与属性链，
\textbf{砍掉} \(\forall\)、\(\sqcup\)、否定、基数限制与逆属性。
被砍掉的恰恰是 5.2 节里制造分支与传播的元凶，
于是 EL 内的分类问题从指数级降到\textbf{多项式时间}。
代价换来的速度是惊人的：医学本体 SNOMED CT 有三十多万个类，
Tableau 推理器分类一次以小时计，ELK 在普通笔记本上以秒计。

工程启示有两条。
其一，\textbf{表达力是花钱的}：如果你的本体 95\% 的公理只是层次加存在约束，
不要为剩下 5\% 的华丽构造把整库拖进 Tableau 的指数世界——
要么改写掉那 5\%，要么把它们隔离到单独的小模块（5.7 节的模块化）。
其二，OWL 2 的三个 Profile（EL、QL、RL）本质上是三份「预谈好价格的合同」：
EL 押注大规模类层次（生物医学、产品分类），
QL 押注查询重写（把推理编译进 SQL/SPARQL，5.1.6 的后向风格），
RL 押注规则物化（下一小节的前向风格）。
选 Profile 就是选推理算法，第 4 章语法层面的 Profile 介绍，在这里有了算法层面的注脚。

取舍长什么样，拿本章自己的公理试一下就知道。
设备类层次，以及「数控车床必有数控系统」那样的存在约束（5.2.6 第一条），EL 全盘照收；
而 5.2.6 那条用 \(\forall\) 锁利用率区间的公理，EL 写不了。
怎么办？三选一：把它改写成数据属性的值域声明（多数场景语义足够）；
把校验意图挪去 SPARQL／SHACL 这类闭世界检查（它本来就更像数据校验，见 5.2.6 的提醒）；
或者把少数离不开 \(\forall\) 的公理隔离成小模块，单独交给 Tableau 推理器、低频运行。
\textbf{Profile 取舍不是「忍痛割爱」，而是把每条公理送到对它最便宜的执行引擎去}。

\subsection{规则引擎与 RETE 网络}

SWRL（5.4 节）这类规则一旦多起来，朴素执行方式撑不住：
每来一条新事实就把所有规则的所有前件从头匹配一遍，规则数 \(\times\) 事实数的乘积很快爆炸。
产生式系统社区四十年前给出的答案沿用至今：\textbf{RETE 算法}。

RETE 的直觉是\textbf{把规则编译成一张筛网，让数据流过网，而不是拿规则反复扫数据}。
网络分两段：\textbf{alpha 段}做单条件初筛——
「是设备吗」「功率大于 10 吗」，每个原子条件一个筛孔；
\textbf{beta 段}做跨条件拼接——「这台设备」与「这条功率记录」说的是不是同一个 \texttt{?x}，
拼接成功的半成品（部分匹配）被\textbf{缓存}在节点里。
一条新事实到达时，只沿着它能通过的筛孔向下渗透，
与各节点缓存的半成品增量拼接，到达网络底部的事实组合即触发规则开火。
\textbf{用空间换时间、用增量换全量}是其全部精髓：
旧事实的匹配成果留在缓存里永不重算，
代价是内存里长期挂着大量半成品——规则越多、连接变量越多，内存越吃紧。

用 5.4 的高功率规则走一遍数据流。
规则「设备(\texttt{?x}) \(\land\) 功率(\texttt{?x}, \texttt{?p}) \(\land\)
\texttt{?p} \(>\) 10 \(\rightarrow\) 高功率设备(\texttt{?x})」编译成网络后：
alpha 段三个筛孔——「是设备断言吗」「是功率断言吗」「数值大于 10 吗」；
beta 段一个连接节点，左缓存存「已知的设备」，右缓存存「已过阈值的功率记录」。
事实「设备(Lathe\_003)」先到，穿过第一孔，存入左缓存，无可拼接，流程安静结束；
随后「功率(Lathe\_003, 15)」到达，连过两孔进入右缓存，
与左缓存中的 Lathe\_003 按变量 \texttt{?x} 拼接成功——规则开火。
若先到的是另一台 Mill\_009 的功率记录而它尚未被声明为设备，
这条半成品就留在右缓存里\textbf{等}，直到设备断言补到为止。
「数据到一条、网里渗一条、缓存接住没凑齐的」——这就是 RETE 的全部日常。
Drools、Jena 前向引擎以及多数商用规则引擎都是 RETE 或其变体；
Semantica 的受限正向规则 API 也体现“事实驱动规则”的基本节奏，但不因此宣称采用
完整 RETE、支持一般 SWRL 或拥有第三方引擎的全部 profile。

\subsection{物化还是查询时推理：一张架构对照表}

前向、后向两种策略落到系统架构上，就是\textbf{物化}（materialization，
把推理结论预先算出并写回三元组库）与\textbf{查询时推理}
（query-time reasoning，查询到来才推导）的取舍。
这是搭知识图谱平台绕不开的第一个架构决策，先看对照表再谈选择。

\begin{center}
\begin{tabularx}{\linewidth}{@{}p{0.20\linewidth}XX@{}}
\toprule
\textbf{维度} & \textbf{物化（写时推理）} & \textbf{查询时推理（读时推理）} \\
\midrule
推理成本付在 & 数据写入/批处理阶段 & 每次查询执行阶段 \\
查询延迟 & 最低，等同普通图查询 & 较高，且随推理深度波动 \\
存储开销 & 膨胀明显（隐含三元组全部落库） & 无额外存储 \\
数据更新 & 代价高：删改需维护推理闭包 & 代价低：下次查询自然反映 \\
结论新鲜度 & 取决于物化批次，可能滞后 & 永远最新 \\
典型实现 & OWL RL 规则物化、RETE 引擎 & OWL QL 查询重写、后向规则 \\
适合场景 & 读多写少、低延迟看板、批量分析 & 写多读少、数据高频变动、即席诊断 \\
\bottomrule
\end{tabularx}
\end{center}

经验法则：\textbf{更新频率与查询频率之比}是第一判据。
设备台账、产品分类这类「慢知识」读多写少，物化几乎总是赢家；
传感器状态这类「快知识」每秒都变，物化闭包维护不过来，留给查询时处理或干脆不进推理层。
成熟平台普遍走\textbf{分层混合}：稳定的 TBox 层次物化进库，
快变的 ABox 走轻量查询时规则——怎么把「物化哪些、何时重算」做成可运维的管道，
是第 7 章知识图谱管理的主题之一，5.7.4 会先给出衔接的原则。

\section{SWRL 深入：把业务规则写进本体}

描述逻辑擅长概念结构，却写不了「如果这台设备正在运行\textbf{并且}有任务派给它，
\textbf{那么}产生一个冲突告警」这种横跨多个个体、带比较运算的业务规则。
\textbf{SWRL}（Semantic Web Rule Language）补上这一环：
它让 Horn 风格的规则与 OWL 本体共用同一套类与属性词汇表，
规则推出的结论直接成为本体事实，继续参与 OWL 推理。
本节先解剖语法，再讲保住可判定性的 DL-Safe 限制，
然后把仓库规则库的五类调度规则逐类精讲，最后划清 SWRL 做不到的事。

\subsection{一条规则的解剖：三类原子}

SWRL 规则的形态是「前件 \(\rightarrow\) 后件」：
前件（antecedent）是用 \(\land\) 连接的一串\textbf{原子}（atom），全部成立时，
后件（consequent）的原子被断言为新事实。变量写作 \texttt{?x} 这样的问号标识符。
原子只有寥寥几类，却足以覆盖大部分业务规则：

\begin{itemize}
\item \textbf{类原子}（class atom）：形如「高功率设备(\texttt{?x})」，断言个体属于某类。
  注意\textbf{类原子有且只有一个参数}——
  想表达「两个任务时间冲突」不能写成 TimeConflict(\texttt{?t1}, \texttt{?t2})，
  只能给其中一方贴单参数标签，这个细节在下文时间冲突规则里会再次出现。
\item \textbf{属性原子}（property atom）：形如「功率(\texttt{?x}, \texttt{?p})」，恰好两个参数，
  对象属性连接两个个体，数据属性连接个体与字面值。
\item \textbf{内建原子}（built-in atom）：形如 \texttt{swrlb:greaterThan(?p, 10)}，
  调用内建函数库做比较、算术与字符串处理，是规则拥有「计算能力」的来源。
\item 另有同一性原子 sameAs / differentFrom，用于声明两个变量绑定的个体相同或不同，工程中出场较少。
\end{itemize}

包内教学资产 \texttt{swrl-rules} 末尾的「语法备注」段落把这些约定浓缩成了速查注释，
读规则前先扫一眼，可以避开「类原子塞两个参数」这类最常见的语法错误。

\subsection{swrlb 内建函数族}

内建原子的命名空间前缀 \texttt{swrlb:} 背后是一整族函数，按用途分三组，常用成员如下。

\begin{center}
\begin{tabularx}{\linewidth}{@{}p{0.14\linewidth}p{0.44\linewidth}X@{}}
\toprule
\textbf{函数族} & \textbf{常用成员} & \textbf{说明与示例} \\
\midrule
比较族 & \texttt{equal} / \texttt{notEqual} / \texttt{lessThan} / \texttt{lessThanOrEqual} / \texttt{greaterThan} / \texttt{greaterThanOrEqual} & 数值与可比类型的大小判断；「功率大于 10」即 \texttt{swrlb:greaterThan(?p, 10)} \\
算术族 & \texttt{add} / \texttt{subtract} / \texttt{multiply} / \texttt{divide} / \texttt{mod} / \texttt{abs} & 第一个参数是结果变量：\texttt{swrlb:subtract(?d, ?end, ?start)} 即 \texttt{?d = ?end - ?start} \\
字符串族 & \texttt{stringConcat} / \texttt{substring} / \texttt{stringLength} / \texttt{upperCase} / \texttt{lowerCase} / \texttt{contains} / \texttt{matches} & 编号解析、前缀匹配；如用 \texttt{matches} 校验设备编号格式 \\
\bottomrule
\end{tabularx}
\end{center}

两个工程提醒。
其一，内建原子是\textbf{纯计算谓词}，只能对\textbf{已绑定}的变量求值，不会反向生成绑定——
\texttt{swrlb:subtract(?d, ?end, ?start)} 要求 \texttt{?end} 与 \texttt{?start}
已经被前面的属性原子绑定，否则规则永不触发；写规则时习惯把属性原子放前、内建原子放后。
其二，内建函数没有「当前时间」这类\textbf{有副作用的来源}，这个看似小事的缺口在 5.5 节会变成主角。
顺手给字符串族一个用例：设备编号自带批次语义时，规则
「Equipment(\texttt{?e}) \(\land\) hasSerial(\texttt{?e}, \texttt{?s}) \(\land\)
\texttt{swrlb:contains(?s, "EQ-2023")} \(\rightarrow\) Batch2023(\texttt{?e})」
把编号里的批次前缀翻译成类标签——
老系统的编码规则借一条规则升格为本体知识，数据迁移省掉一列人工映射。

\subsection{DL-Safe：把规则拴在已命名个体上}

OWL DL 之所以敢承诺推理可终止，靠的是对表达力的精细克制；
而把不受限的 Horn 规则与 OWL 拼在一起，组合语言的推理问题是\textbf{不可判定}的——
规则可以与 \(\exists\) 公理联手，在匿名个体的无限链条上无休止地互相触发，
5.2.5 的阻塞机制对「规则 + 公理」的混合体系不再够用。

工程上的解药是 \textbf{DL-Safe 限制}：\textbf{规则变量只允许绑定到 ABox 中显式命名的个体}，
推理器虚构的匿名个体（比如 5.2 节的 cooling\_1）对规则不可见。
无限链条的来源被掐断——命名个体是有限的，规则的触发组合就是有限的，可判定性失而复得。
代价是完备性的让步：规则「看不见」纯逻辑世界里的匿名证人，
个别本该触发的结论会漏掉。
实践中这几乎不构成困扰——业务规则关心的本来就是真实台账里那些有名有姓的设备与任务；
主流实现（Pellet 的 SWRL 支持、HermiT 等）默认就在 DL-Safe 语义下执行规则。
记住一句话即可：\textbf{SWRL 在工程中的标准用法就是 DL-Safe 用法，规则管命名个体，公理管概念世界}。

「漏掉的结论」长什么样？接上 5.2.3 的尾巴。
推演造出的匿名证人 cooling\_1 对规则不可见，
于是规则「冷却系统(\texttt{?c}) \(\rightarrow\) 需巡检(\texttt{?c})」
会对台账里命名登记的 CoolSys\_A 触发，对 cooling\_1 沉默。
这是缺陷吗？恰恰相反：巡检清单本来就只该列\textbf{真实存在的资产}；
如果业务上确实有一台冷却系统需要巡检，
正确做法是把它作为命名个体录入台账，而不是指望规则去巡检一个逻辑脚手架。
DL-Safe 的边界与工程的台账边界天然重合——这是它在实践中几乎无感的真正原因。

\subsection{五类调度规则逐条精讲}

仓库规则库覆盖车间调度的五类场景，下面逐类过：每类先讲业务动机，再讲推理链怎么走。

\textbf{第一类：基础推理规则——属性阈值触发分类。}
业务动机：维护手册规定功率超过 10 千瓦的设备必须配冷却，
调度系统希望这个判断自动完成，而不是靠人工维护「高功率设备清单」。

\begin{codebox}
\codeline{\textcolor{cmtgray}{\#\ 规则：高功率设备需要冷却}}
\codeline{\textcolor{cmtgray}{\#\ Rule:\ High-power\ equipment\ requires\ cooling}}
\codeline{设备(?x)\ \(\land\)\ 功率(?x,\ ?p)\ \(\land\)\ swrlb:greaterThan(?p,\ 10)}
\codeline{\hspace*{4\ttcw}\(\rightarrow\)\ 高功率设备(?x)}
\codeblank{}
\codeline{高功率设备(?x)\ \(\rightarrow\)\ 需要冷却(?x)}
\codeblank{}
\codeline{\textcolor{cmtgray}{\#\ 验证推理链（前向推理示例）：}}
\codeline{\textcolor{cmtgray}{\#\ 事实:\ 设备(Lathe\_003)，功率(Lathe\_003,\ 15)}}
\codeline{\textcolor{cmtgray}{\#\ 步骤1:\ 功率(Lathe\_003,\ 15)\ \(\land\)\ 15>10\ \(\rightarrow\)\ 高功率设备(Lathe\_003)}}
\codeline{\textcolor{cmtgray}{\#\ 步骤2:\ 高功率设备(Lathe\_003)\ \(\rightarrow\)\ 需要冷却(Lathe\_003)}}
\codeline{\textcolor{cmtgray}{\#\ 结论:\ Lathe\_003\ 需要冷却}}
\end{codebox}
\color{black}

\assetsource{ch05}{swrl-rules}

推理链是两级火箭：第一条规则用属性原子绑定功率值，内建原子做阈值比较，
满足即给设备贴上「高功率设备」类标签；
第二条规则接力，把类标签翻译成「需要冷却」的结论。
片段下方的验证步骤值得手过一遍：事实「功率(Lathe\_003, 15)」进来，
15 \(>\) 10 成立，步骤 1 推出高功率设备(Lathe\_003)，步骤 2 推出需要冷却(Lathe\_003)——
一条数据触发两跳推理，这正是 5.1.6 所说前向推理的微缩样本。
还可以与 5.2.3 对照：同一个业务结论，DL 公理版靠 Tableau 的等价类触发与 \(\exists\) 造点，
SWRL 版靠两条规则接力——\textbf{概念建模与规则建模是同一知识的两种落地姿势}，
取舍标准在 5.4.5 讨论。

\textbf{第二类：设备调度规则——状态、时间与维护的三重冲突检测。}
业务动机：派工最怕三件事——把任务派给正在运行的设备、
两个任务在同一台设备上撞档期、设备维护超时无人过问。三条规则分别站岗。

\begin{codebox}
\codeline{\textcolor{cmtgray}{\#\ ==============================}}
\codeline{\textcolor{cmtgray}{\#\ 2.\ 设备调度规则}}
\codeline{\textcolor{cmtgray}{\#\ ==============================}}
\codeblank{}
\codeline{\textcolor{cmtgray}{\#\ 规则：状态冲突检查\ -\ 正在运行的设备不能分配新任务}}
\codeline{\textcolor{cmtgray}{\#\ Rule:\ Status\ conflict\ -\ running\ equipment\ cannot\ be\ assigned\ new\ tasks}}
\codeline{Equipment(?e)\ \(\land\)\ hasStatus(?e,\ Running)\ \(\land\)}
\codeline{Task(?t)\ \(\land\)\ assignedTo(?t,\ ?e)}
\codeline{\hspace*{4\ttcw}\(\rightarrow\)\ StatusConflict(?t)}
\codeblank{}
\codeline{\textcolor{cmtgray}{\#\ 规则：时间冲突检查}}
\codeline{\textcolor{cmtgray}{\#\ Rule:\ Time\ conflict\ detection}}
\codeline{\textcolor{cmtgray}{\#\ 注：后件使用单变量类断言\ TimeConflict(?t1)，SWRL类原子只接受单个参数}}
\codeline{\textcolor{cmtgray}{\#\ Note:\ The\ consequent\ uses\ a\ unary\ class\ assertion\ TimeConflict(?t1);}}
\codeline{\textcolor{cmtgray}{\#\ \ \ \ \ \ \ SWRL\ class\ atoms\ take\ exactly\ one\ argument.}}
\codeline{Task(?t1)\ \(\land\)\ Task(?t2)\ \(\land\)}
\codeline{assignedTo(?t1,\ ?e)\ \(\land\)\ assignedTo(?t2,\ ?e)\ \(\land\)}
\codeline{hasStartTime(?t1,\ ?s1)\ \(\land\)\ hasEndTime(?t1,\ ?e1)\ \(\land\)}
\codeline{hasStartTime(?t2,\ ?s2)\ \(\land\)\ hasEndTime(?t2,\ ?e2)\ \(\land\)}
\codeline{swrlb:lessThan(?s1,\ ?e2)\ \(\land\)}
\codeline{swrlb:greaterThan(?e1,\ ?s2)}
\codeline{\hspace*{4\ttcw}\(\rightarrow\)\ TimeConflict(?t1)}
\codeblank{}
\codeline{\textcolor{cmtgray}{\#\ 规则：长时间维护预警（超过24小时）}}
\codeline{\textcolor{cmtgray}{\#\ Rule:\ Long\ maintenance\ warning\ (over\ 24\ hours)}}
\codeline{\textcolor{cmtgray}{\#\ 注：swrlb:add(?now,\ 86400,\ ?tomorrow)\ 中\ ?now}}
\codeline{\hspace*{4\ttcw}需由外部提供或使用swrlb:currentTimeMillis等内置函数}
\codeline{\textcolor{cmtgray}{\#\ Note:\ The\ built-in\ swrlb:add\ requires\ ?now\ to\ be\ bound.\ Use\ a\ dedicated\ temporal}}
\codeline{\hspace*{4\ttcw}built-in}
\codeline{\textcolor{cmtgray}{\#\ \ \ \ \ \ \ or\ inject\ the\ current\ timestamp\ as\ an\ ABox\ fact\ before\ reasoning.}}
\codeline{Equipment(?e)\ \(\land\)\ hasStatus(?e,\ Maintenance)\ \(\land\)}
\codeline{hasMaintenanceStart(?e,\ ?start)\ \(\land\)}
\codeline{hasMaintenanceEnd(?e,\ ?end)\ \(\land\)}
\codeline{swrlb:subtract(?end,\ ?start,\ ?duration)\ \(\land\)}
\codeline{swrlb:greaterThan(?duration,\ 86400)}
\codeline{\hspace*{4\ttcw}\(\rightarrow\)\ LongMaintenanceWarning(?e)}
\codeblank{}
\end{codebox}
\color{black}

\assetsource{ch05}{swrl-rules}

状态冲突规则的推理链最直白：设备状态为 Running，又有任务 \texttt{?t} 派给它，
两个条件在 \texttt{?e} 上会合，立刻给 \texttt{?t} 贴 StatusConflict 标签。
注意它检查的是「已派单」之后的矛盾——SWRL 只能断言新事实，不能阻止旧事实写入，
所以工程形态是「写入后检测 + 告警回滚」而非「写入前拦截」。

时间冲突规则是五类中最值得逐字读的一条。
两个任务各自带起止时间 [\texttt{?s1}, \texttt{?e1}] 与 [\texttt{?s2}, \texttt{?e2}]，
重叠判据浓缩成两个内建原子：\texttt{?s1} \(<\) \texttt{?e2} 且 \texttt{?e1} \(>\) \texttt{?s2}。
不妨代入数字验证：任务一占 8 点到 10 点，任务二占 9 点到 11 点，
8 \(<\) 11 且 10 \(>\) 9，两式同时成立，冲突命中；
若任务二改到 10 点到 12 点（首尾相接不算重叠），第二式 10 \(>\) 10 不成立，规则安静放行。
这个「双不等式」其实是 5.5.4 节 Allen 区间代数的工程速记，到时再正式对账。
另外注意后件写的是单参数的 TimeConflict(\texttt{?t1})——正是 5.4.1 强调的类原子限制；
想记录「与谁冲突」，可以在后件追加一个属性原子 conflictsWith(\texttt{?t1}, \texttt{?t2})。

长时间维护预警规则展示算术内建的标准用法：
结束时间减开始时间得到时长，超过 86400 秒（恰好 24 小时）触发预警。
片段注释里藏着一个重要伏笔：规则里若想引用「现在」，
变量必须由外部注入——SWRL 没有时钟，5.5.3 专门解决这个问题。

\textbf{第三类：工艺推荐规则——用相似性搬运经验。}
业务动机：新产品来了，工艺员的第一反应是「找一个材料相同、特征相同的老产品，抄它的工艺」。
规则把这个动作自动化：
\texttt{?p1} 与 \texttt{?p2} 同材料、同特征，且 \texttt{?p2} 已有工艺 \texttt{?proc}，
则推出 recommendedProcess(\texttt{?p1}, \texttt{?proc})。
推理链的本质是\textbf{沿共同属性做个体间的知识迁移}，
配套的第二条规则再按工序类型与设备能力推荐承接设备（车削工序配有车削能力的车床）。
注意推荐结论用的是专门属性 recommendedProcess 而非直接写 hasProcess——
推理产物与人工确认的事实分属性存放，这是第 7 章数据治理会反复强调的「推理结果可追溯」原则。
此类两条规则的全文见包内教学资产第 3 节。

\textbf{第四类：工序顺序规则——把工艺纪律写成逻辑。}
业务动机：工序有先后，错序是质量事故的温床——焊缝没做热处理就交付，残余应力迟早出问题。

\begin{codebox}
\codeline{\textcolor{cmtgray}{\#\ ==============================}}
\codeline{\textcolor{cmtgray}{\#\ 4.\ 工序顺序推理规则}}
\codeline{\textcolor{cmtgray}{\#\ ==============================}}
\codeblank{}
\codeline{\textcolor{cmtgray}{\#\ 规则：工序必须按顺序执行（前序完成才能开始后序）}}
\codeline{\textcolor{cmtgray}{\#\ Rule:\ Operations\ must\ follow\ order\ (predecessor\ must\ complete\ before\ successor}}
\codeline{\hspace*{4\ttcw}starts)}
\codeline{Operation(?op1)\ \(\land\)\ Operation(?op2)\ \(\land\)}
\codeline{precedes(?op1,\ ?op2)\ \(\land\)\ hasStatus(?op1,\ Completed)}
\codeline{\hspace*{4\ttcw}\(\rightarrow\)\ canStart(?op2,\ true)}
\codeblank{}
\codeline{\textcolor{cmtgray}{\#\ 规则：热处理必须在焊接之后进行（领域约束）}}
\codeline{\textcolor{cmtgray}{\#\ Rule:\ Heat\ treatment\ must\ follow\ welding\ (domain\ constraint)}}
\codeline{Operation(?op1)\ \(\land\)\ Operation(?op2)\ \(\land\)}
\codeline{hasType(?op1,\ HeatTreatment)\ \(\land\)\ hasType(?op2,\ Welding)\ \(\land\)}
\codeline{isPartOfSameProcess(?op1,\ ?op2)}
\codeline{\hspace*{4\ttcw}\(\rightarrow\)\ precedes(?op2,\ ?op1)}
\codeblank{}
\codeline{\textcolor{cmtgray}{\#\ 规则：工艺参数超范围预警}}
\codeline{\textcolor{cmtgray}{\#\ Rule:\ Process\ parameter\ out-of-range\ warning}}
\codeline{CuttingOperation(?op)\ \(\land\)\ hasSpeed(?op,\ ?s)\ \(\land\)}
\codeline{swrlb:greaterThan(?s,\ 5000)}
\codeline{\hspace*{4\ttcw}\(\rightarrow\)\ ParameterWarning(?op)}
\codeblank{}
\end{codebox}
\color{black}

\assetsource{ch05}{swrl-rules}

第一条规则维护执行秩序：前序工序状态变为 Completed，
后序工序才获得 canStart 的放行标记，调度器只对放行的工序派工。
第二条规则更有味道：它的后件不是告警，而是\textbf{推出一条 precedes 关系}——
「同一工艺流程内，热处理必须排在焊接之后」这条领域纪律，
被翻译成「检测到这两类工序共处一个流程，就自动补上先后约束」。
新补的 precedes 事实又会成为第一条规则的前件输入——\textbf{规则之间通过共享词汇表自动接力}，
这是把规则写进本体（而不是写进应用代码）的最大红利：组合行为不需要任何人显式编排。
第三条规则用 5000 转的阈值给切削参数站岗，结构与第一类的阈值规则同型。

\textbf{第五类：资源冲突规则——可用量见底就报警。}
业务动机：任务声明需要某资源，而资源的可用量 \texttt{?a} 已小于 1（即无可分配单位），
则给任务贴 ResourceConflict 标签。
推理链与状态冲突同构，差别只在比较对象从状态枚举换成了库存数值；
规则全文见包内 \texttt{swrl-rules} 教学资产第 5 节。
有一个分工细节值得点破：规则只\textbf{读} hasAvailability 的数值并断言冲突，
\textbf{不做扣减}——库存增减是库存系统的事务职责，
推理层保持「只产出结论、不修改输入」的纯净性，
否则规则反复触发时连自己消费过哪个版本的库存都说不清。
把五类规则连起来看就是一份\textbf{派工前置检查清单}：
状态、档期、维护时长、工序秩序、资源余量各设一道闸，
任何一道闸亮灯，调度器都能在派单生效前拿到带标签的告警事实。
五类规则合在一起，构成第 9 章规则需求的来源规格；但当前 ch09 只结构化保留
这些 SWRL，不执行一般 SWRL/built-ins。可复算的只有 ch05 包中声明的受限正向链，
不得用旧 Python 实现填补 runtime 的红灯。

\subsection{SWRL 写不了什么：否定的边界与绕法}

规则库文件开头有一段加粗的注意事项，值得原文转述其要点：
\textbf{标准 SWRL 不支持规则体中的否定}。
「设备\textbf{不能}加工某材料则告警」「任务\textbf{没有}分配设备则提示」——
这类以「不、没有」开头的条件在 SWRL 前件里没有合法写法。
根源还是第 2 章的开放世界假设：
「知识库里查不到 canProcess(\texttt{?e}, 钛合金)」不等于「这台设备不能加工钛合金」，
缺失只是未知；在 OWA 之上引入「查不到即为假」的否定（negation as failure），
语义与 OWL 的单调推理直接打架，标准因此干脆不提供。

工程上有三条绕法，按侵入性从低到高排列：
\textbf{其一，显式否定属性}——建模时就把「不能加工」做成正面属性 cannotProcess，
用数据把否定说出口，规则照常写正原子；代价是数据维护翻倍。
\textbf{其二，应用层 SPARQL 兜底}——把含否定的判断挪出规则层，
用 \texttt{FILTER NOT EXISTS} 在查询时做局部封闭世界检查；
这是仓库注释推荐的主路线，也是实践中最常见的分工：
\textbf{规则层管正向推导，查询层管缺失检测}。
\textbf{其三，换引擎}——SPARQL Inferencing Notation（SPIN/SHACL 规则）或
支持分层否定的 Datalog 方言允许受控的否定，
代价是离开 SWRL 标准生态。
没有免费午餐：选哪条路，取决于团队更愿意维护数据、维护查询还是维护第二套引擎。

\section{时态知识：给三元组装上时间轴}

到目前为止，知识库里的世界是静止的：Lathe\_003 位于 WorkCell\_B——永远位于。
真实车间却一刻不停：设备搬迁、状态翻转、任务起止。
本节解决「知识如何随时间变化」：先看 RDF 为什么天生没有时间维度，
再给出工程主流方案具名图，然后补上时态规则与 Allen 区间代数这两件推理武器。

\subsection{三元组为什么没有时间}

RDF 的原子单位是三元组：主语—谓语—宾语，
「:Lathe\_003 :locatedIn :WorkCell\_B」。
这是一个\textbf{二元关系}的断言，语法上没有给「何时成立」留任何位置。
想说「3 月 21 日位于 B 区、22 日搬到 A 区」，朴素写法是把两条 locatedIn 都写进库——
但这两条三元组没有任何标记区分新旧，查询会同时返回两个位置，知识库变成一锅时间大杂烩。

经典补救是\textbf{具体化}（reification）或 n 元关系模式：
把「位于」升格为一个对象，再挂上时间属性。
具体看一眼代价：原本一条「:Lathe\_003 :locatedIn :WorkCell\_B」，
具体化后变成一个位置事件对象 \texttt{:loc\_0321}，
身上挂着主体（\texttt{:subject :Lathe\_003}）、地点（\texttt{:place :WorkCell\_B}）、
起始与结束时间（\texttt{:validFrom} / \texttt{:validTo}）——
一条事实膨胀成五条三元组；
查「它现在在哪」也从一跳变成「找事件对象、连三跳、再按时间过滤」的长查询。
第 4 章讲过 n 元模式在「关系本身带属性」时是正解，
但若\textbf{每一条}状态事实都要这么包装，建模噪音就压过了语义收益。
RDF 生态最终走通的路线是换一个维度：\textbf{不改造三元组，而是给三元组分组}。

\subsection{具名图：给状态拍快照}

\textbf{具名图}（Named Graph）把三元组库从「一大袋三元组」升级为「多个有名字的袋子」：
每个图有自己的 URI 作为图名，三元组归属于图，
「三元组 + 图名」合成四元组（quad）。
把图名当作时间戳使用，每个图就是一张\textbf{状态快照}：

\begin{codebox}
\codeline{\textcolor{cmtgray}{\#\ ==============================}}
\codeline{\textcolor{cmtgray}{\#\ 具名图（Named\ Graph）示例}}
\codeline{\textcolor{cmtgray}{\#\ ==============================}}
\codeline{\textcolor{cmtgray}{\#\ Named\ Graph\ 允许对\ RDF\ 三元组进行分组和命名}}
\codeline{\textcolor{cmtgray}{\#\ 常用于表达时间快照或来源元数据}}
\codeblank{}
\codeline{\textcolor{cmtgray}{\#\ TriG\ 格式（Turtle的扩展，支持具名图）：}}
\codeline{\textcolor{cmtgray}{\#\ 具名图的图名必须是URI（用尖括号括起来），不能是裸名}}
\codeline{\textcolor{cmtgray}{\#\ Named\ graph\ names\ must\ be\ URIs\ enclosed\ in\ angle\ brackets,\ not\ bare\ identifiers.}}
\codeblank{}
\codeline{@prefix\ :\ <http://example.org/manufacturing\#>\ .}
\codeline{@prefix\ xsd:\ <http://www.w3.org/2001/XMLSchema\#>\ .}
\codeblank{}
\codeline{\textcolor{cmtgray}{\#\ 2024年3月21日的状态快照}}
\codeline{<http://example.org/graph/20240321>\ \{}
\codeline{\hspace*{4\ttcw}:Lathe\_003\ :status\ "Idle"\ ;\ \ \ \ \ \ \ \ \ \ \ \#\ 车床状态为闲置}
\codeline{\hspace*{15\ttcw}:locatedIn\ :WorkCell\_B\ .\ \ \ \#\ 位于工作单元B}
\codeline{\}}
\codeblank{}
\codeline{\textcolor{cmtgray}{\#\ 2024年3月22日的状态快照（设备已被移动）}}
\codeline{<http://example.org/graph/20240322>\ \{}
\codeline{\hspace*{4\ttcw}:Lathe\_003\ :status\ "Running"\ ;}
\codeline{\hspace*{15\ttcw}:locatedIn\ :WorkCell\_A\ .\ \ \ \#\ 已移动到A区}
\codeline{\}}
\codeblank{}
\end{codebox}
\color{black}

\assetsource{ch05}{temporal-reasoning}

片段用 TriG 格式（Turtle 的具名图扩展）记录了两天的车间状态：
21 日的图里 Lathe\_003 闲置在 B 区，22 日的图里它已在 A 区运行。
两组「互相矛盾」的断言因为住在不同的袋子里而和平共处；
SPARQL 用 \texttt{GRAPH} 子句按图名取数，即可做\textbf{时间切片回溯}——
「21 日盘点时这台车床在哪」一查便知。
两个语法细节别错过：图名必须是 URI（尖括号括起，不能用裸名），
这样图名本身也能成为三元组的主语——
给图挂上「快照时刻、数据来源、采集系统」等元数据，时间戳就升级成了完整的溯源记录。
另外快照粒度是设计决策：按天、按班次还是按事件触发，
取决于回溯需求与存储预算。
快照还支持\textbf{跨期对比}：同一条查询分别套上两个图名执行，
结果做差即是「这两天之间发生了什么变化」——
本例中差集恰好是 Lathe\_003 的状态翻转与搬迁两条，
设备异动报表由此免去专门的变更日志表。
具名图同时也是第 7 章存储策略的主角——
四元组存储（quad store）、按图分区、按图做访问控制，都建立在本节这套机制之上。

\subsection{时态规则与 NOW() 难题}

快照解决「存」，推理还要解决「算」。两条时态规则展示了典型模式与典型陷阱。

\begin{codebox}
\codeline{\textcolor{cmtgray}{\#\ ==============================}}
\codeline{\textcolor{cmtgray}{\#\ 时态推理规则}}
\codeline{\textcolor{cmtgray}{\#\ ==============================}}
\codeblank{}
\codeline{\textcolor{cmtgray}{\#\ 推理规则：在某时间段内运行的设备，该时段内不可用于其他任务}}
\codeline{\textcolor{cmtgray}{\#\ 注：以下为概念性SWRL表示，?t1/?t2\ 为时间段端点，需以xsd:dateTime数据值绑定}}
\codeline{Device(?d)\ \(\land\)\ hasStatus(?d,\ Running,\ ?t1,\ ?t2)\ \(\land\)}
\codeline{Task(?task)\ \(\land\)\ hasTimeInterval(?task,\ ?s,\ ?e)\ \(\land\)}
\codeline{swrlb:greaterThan(?e,\ ?t1)\ \(\land\)\ swrlb:lessThan(?s,\ ?t2)\ \(\land\)}
\codeline{swrlb:notEqual(?task,\ ?currentTask)}
\codeline{\hspace*{4\ttcw}\(\rightarrow\)\ unavailable(?d,\ ?task)}
\codeblank{}
\codeline{\textcolor{cmtgray}{\#\ 推理规则：设备历史故障会影响维护计划}}
\codeline{\textcolor{cmtgray}{\#\ 注：SWRL没有内置NOW()函数；实际应用中需在规则执行前将当前时间注入ABox，}}
\codeline{\textcolor{cmtgray}{\#\ \ \ \ \ 或使用支持时态推理的规则引擎扩展。}}
\codeline{\textcolor{cmtgray}{\#\ Note:\ SWRL\ has\ no\ NOW()\ built-in.\ Inject\ the\ current\ timestamp\ as\ an\ ABox\ fact}}
\codeline{\textcolor{cmtgray}{\#\ \ \ \ \ \ \ (e.g.,\ :currentTime\ :value\ "2024-03-22T08:00:00"\^{}\^{}xsd:dateTime)\ before}}
\codeline{\hspace*{4\ttcw}reasoning.}
\codeline{Device(?d)\ \(\land\)\ hasFaultRecord(?d,\ ?fault)\ \(\land\)}
\codeline{hasFaultTime(?fault,\ ?t)\ \(\land\)}
\codeline{currentTime(:system,\ ?now)\ \(\land\)}
\codeline{swrlb:subtract(?now,\ ?t,\ ?elapsed)\ \(\land\)}
\codeline{swrlb:lessThan(?elapsed,\ 2592000)}
\codeline{\hspace*{4\ttcw}\(\rightarrow\)\ requiresInspection(?d,\ true)}
\end{codebox}
\color{black}

\assetsource{ch05}{temporal-reasoning}

第一条规则是 5.4.4 时间冲突规则的时态版：
设备在 [\texttt{?t1}, \texttt{?t2}] 区间运行，任务区间 [\texttt{?s}, \texttt{?e}] 与之相交
（同样的双不等式判据），则推出设备对该任务不可用。
注释提醒这是\textbf{概念性写法}：标准 SWRL 的属性原子只有两个参数，
hasStatus(\texttt{?d}, Running, \texttt{?t1}, \texttt{?t2}) 这种四参形态实际要靠
n 元模式落地——状态本身建为对象，时间端点挂在状态对象上。

第二条规则暴露一个更根本的缺口：「故障发生在最近 30 天内」需要\textbf{当前时间}，
而 \textbf{SWRL 没有 NOW() 内建}——规则语言是纯声明式的，没有时钟这种副作用源。
工程绕法已经写在规则里：把当前时间作为一条普通事实
currentTime(:system, \texttt{?now}) 注入 ABox，规则像读数据一样读「现在」。
配套的运维纪律是\textbf{每轮推理前刷新这条事实}（由调度脚本写入），
推理结果的「新鲜度」就由刷新频率决定。
顺手验算规则里的常量：2592000 秒 \(=\) 30 \(\times\) 86400，恰是 30 天，
与注释宣称的「近一个月故障需要复检」吻合——
规则里的魔法数字都应当这样配上注释并可手工验算。

\subsection{Allen 区间代数十三式}

时间冲突写了两回「双不等式」，是时候给它正名。
两个时间区间之间的相对位置关系，逻辑学家 Allen 证明了\textbf{恰好有十三种}，
彼此互斥且穷尽所有情况——这就是\textbf{Allen 区间代数}（interval algebra），
所有时态推理的坐标系。十三种关系是六对互逆加一个自逆：

\begin{center}
\begin{tabularx}{\linewidth}{@{}p{0.26\linewidth}p{0.20\linewidth}X@{}}
\toprule
\textbf{关系（区间 X 对 Y）} & \textbf{逆关系} & \textbf{直觉含义（X 视角）} \\
\midrule
before（在前） & after（在后） & X 整体结束于 Y 开始之前，中间有空隙 \\
meets（紧邻） & met-by（被紧邻） & X 的终点恰好是 Y 的起点，首尾相接无重叠 \\
overlaps（交叠） & overlapped-by（被交叠） & X 先开始，Y 在 X 结束前开始，两段部分重叠 \\
starts（同始） & started-by（被同始） & 两者同时开始，X 先结束 \\
during（包含于） & contains（包含） & X 完全落在 Y 内部，首尾都不贴边 \\
finishes（同终） & finished-by（被同终） & 两者同时结束，X 后开始 \\
equals（重合） & ——（自逆） & 起止完全相同 \\
\bottomrule
\end{tabularx}
\end{center}

调度工程最常用三种，逐一对账。
\textbf{before} 是工序顺序的语义基座：5.4.4 的 precedes 关系，
严格说断言的就是「前序区间 before（或 meets）后序区间」；
排程算法输出的甘特图，本质是给所有工序对指派 before/meets 关系。
\textbf{during} 是窗口约束的语言：「维护窗口期内禁止派工」即
「任务区间 during 维护区间则冲突」；「夜班时段的任务」即「任务区间 during 夜班区间」。
\textbf{overlaps} 是冲突检测的核心：5.4.4 那条双不等式
\texttt{?s1} \(<\) \texttt{?e2} \(\land\) \texttt{?e1} \(>\) \texttt{?s2}，
对照表格可以精确说出它的覆盖范围——它命中的不只是严格的 overlaps，
而是「两区间存在公共内部」的全部九种关系
（overlaps、overlapped-by、starts、started-by、during、contains、finishes、finished-by、equals），
恰好排除掉无重叠的四种（before、after、meets、met-by）。
工程上的「时间冲突」要的正是这九种的并集——
一条双不等式顶九个关系判断，这就是 Allen 代数给调度工程的速记红利；
反过来，当业务真要区分「部分交叠」与「完全包含」（比如计费规则不同）时，
就必须放弃速记，按表格写出精确的端点比较。

Allen 代数的另一半威力是\textbf{组合推理}（composition）：
已知 X 对 Y、Y 对 Z 的关系，可以推出 X 对 Z 的关系范围。
有些组合结论唯一——「工序 A before B、B before C」必然得出「A before C」，
「A during B、B during C」必然得出「A during C」；
有些组合给出候选集——「A overlaps B、B overlaps C」只能断定
A 对 C 落在 \{before, meets, overlaps\} 三者之一。
排程一致性检测的原理正在于此：把全部工序对的区间关系看作约束网络，
反复用组合表传播、收窄每对工序的候选关系集，
某一对的候选集被挤成空集，就证明这份排程在时间上自相矛盾——
不用真的去算每个时刻，纯靠关系代数就能宣判。

\subsection{W3C Time Ontology 一瞥}

时间建模不必从零造词汇。W3C 的 \textbf{Time Ontology}（命名空间习惯前缀 \texttt{time:}）
把上面这套概念标准化：\texttt{time:Instant}（时刻）与 \texttt{time:Interval}（区间）两个核心类，
\texttt{time:hasBeginning} / \texttt{time:hasEnd} 连接区间与端点时刻，
\texttt{time:inXSDDateTime} 把时刻锚定到具体的 \texttt{xsd:dateTime} 值，
十三种 Allen 关系则原样收录为 \texttt{time:intervalBefore}、\texttt{time:intervalOverlaps}
等对象属性。
此外还有时长与粒度词汇（\texttt{time:hasDuration}、年月日时分秒的分级描述），
足够覆盖排程场景的绝大部分表达需求。
自建调度本体时，把任务起止建成 \texttt{time:Interval} 并复用这些属性，
能换来与排程工具、日历系统的互操作性——标准词汇表的价值在第 4 章已有论述，此处是它在时间维度的具体兑现。
与 5.5.2 的关系一句话说清：具名图负责「这批断言属于哪个时间上下文」，
Time Ontology 负责「时间对象本身怎么描述」，两者一个管袋子、一个管刻度，常常同库并用。

\section{不确定性推理：当知识不再非黑即白}

OWL 的世界是确定性的：一条公理要么成立，要么不成立；
推理器输出的每个结论都板上钉钉。
这种「非真即假」恰恰是工程现场最稀缺的奢侈品——
传感器有噪声，故障是概率事件，「高温的设备\textbf{大概率}会出问题」才是车间知识的真实形态。
本节给确定性骨架补上概率血肉：从给公理标注概率，到一次完整的贝叶斯诊断计算，
再到马尔可夫逻辑网络的权重语义，最后落到「两层分离」的架构原则。

\subsection{确定性知识的边界}

先认清边界在哪里。描述逻辑可以漂亮地表达
「运行中 \(\sqcap\) 维护中 \(\sqsubseteq\) \(\bot\)」——这是\textbf{制度性知识}，违反即矛盾，没有中间地带。
但它无法表达「高功率设备需要定期维护的概率是 85\%」：
OWL 里要么写「高功率设备 \(\sqsubseteq\) 需维护设备」（断言 100\%，太武断），
要么不写（丢失 85\% 这条宝贵经验）。
\textbf{统计性知识在确定性逻辑里没有容身之处}，而预测性维护、质量风险评估、
传感器数据融合，全都靠统计性知识吃饭。
概率本体（probabilistic ontology）的各家方案，本质都是在回答同一个问题：
把概率挂在哪里、怎么跟着推理一起算。

\subsection{给公理标上概率}

最直接的方案是\textbf{在 TBox 公理上标注条件概率}（P-SROIQ、PR-OWL 等风格），
包内教学资产的开头给出两条标准样例：
「P(NeedsMaintenance \(\mid\) HighPowerEquipment) \(=\) 0.85」——
读作「100 台高功率设备中约 85 台需要定期维护」；
「P(PrecisionDrift \(\mid\) LongRunning) \(=\) 0.60」——长时间运行的设备六成会出现精度漂移。
注意概率标注的是\textbf{类包含关系的强度}：
它把第 2 章的 \(\sqsubseteq\)（必然包含）放松成了「以多大比例包含」，
原本二值的箭头有了刻度。

单条概率公理只能做一跳估计；多个不确定因素交织时，需要把它们组织成\textbf{贝叶斯网络}：
节点是随机变量，边表示直接影响，
每个节点带一张\textbf{条件概率表}（CPT，conditional probability table），
给出在父节点各种取值组合下本节点的概率。
仓库的故障诊断网络只有三个节点，却五脏俱全：

\begin{codebox}
\codeline{\textcolor{cmtgray}{\#\ ==============================}}
\codeline{\textcolor{cmtgray}{\#\ 2.\ 贝叶斯网络条件概率表\ (CPT)}}
\codeline{\textcolor{cmtgray}{\#\ ==============================}}
\codeline{\textcolor{cmtgray}{\#\ 设备故障诊断贝叶斯网络结构：}}
\codeline{\textcolor{cmtgray}{\#}}
\codeline{\textcolor{cmtgray}{\#\ \ \ HighTemp（高温）\ \ \ \ \ LongRunning（长时间运行）}}
\codeline{\textcolor{cmtgray}{\#\ \ \ \ \ \ \ \ \ \ \textbackslash{}\ \ \ \ \ \ \ \ \ \ \ \ \ \ \ \ /}}
\codeline{\textcolor{cmtgray}{\#\ \ \ \ \ \ \ \ \ \ \ v\ \ \ \ \ \ \ \ \ \ \ \ \ \ v}}
\codeline{\textcolor{cmtgray}{\#\ \ \ \ \ \ \ \ \ \ \ \ \ Fault（故障）}}
\codeline{\textcolor{cmtgray}{\#}}
\codeline{\textcolor{cmtgray}{\#\ 条件概率表：}}
\codeline{P(Fault\ |\ HighTemp\ \(\land\)\ LongRunning)\ \ \ =\ 0.90}
\codeline{P(Fault\ |\ HighTemp\ \(\land\)\ \(\lnot\)LongRunning)\ \ =\ 0.30}
\codeline{P(Fault\ |\ \(\lnot\)HighTemp\ \(\land\)\ LongRunning)\ \ =\ 0.10}
\codeline{P(Fault\ |\ \(\lnot\)HighTemp\ \(\land\)\ \(\lnot\)LongRunning)\ =\ 0.01}
\codeblank{}
\codeline{\textcolor{cmtgray}{\#\ 先验概率：}}
\codeline{P(HighTemp)\ =\ 0.20\ \ \ \ \ \ \#\ 车间环境中设备高温的先验概率}
\codeline{P(LongRunning)\ =\ 0.35\ \ \ \#\ 设备长时间运行的先验概率}
\codeblank{}
\end{codebox}

\assetsource{ch05}{probabilistic-reasoning}

读表的方法：高温与长时间运行是两个原因节点，共同指向结果节点故障。
CPT 四行覆盖原因的全部四种组合——双重打击下故障概率高达 0.90，
仅高温 0.30，仅长时间运行 0.10，两者皆无时只剩 0.01 的本底故障率。
最下面两行先验概率刻画车间的「常态」：任意时刻约 20\% 的设备处于高温，35\% 在长时间运行。
这张小表浓缩了维修班组多年的经验直觉，而且每个数字都可以用维修台账回归校准——
\textbf{概率本体的知识获取，比逻辑公理更容易与数据闭环}。

数字从哪来？工程上三条来源按可靠性排序。
首选\textbf{台账频率}：「高温且长运行」时段共出现多少次、其中多少次随后记了故障工单，
两数相除就是 CPT 第一行的估计值；
台账不够厚时退而求其次用\textbf{专家校准}——
让老师傅用「一百台里大概几台」的措辞给数，比直接问「概率多少」靠谱得多；
最后，\textbf{小样本要平滑}：某种组合从未出过故障，也不该把概率定成精确的零
（加一平滑之类的技术给「还没见过」留余地），否则一次反例就让模型彻底失语。
还有一笔规模账要算：CPT 的行数随父节点数量\textbf{指数增长}——
k 个二值父节点就是 \(2^{k}\) 行，三个父节点 8 行尚可维护，六个就是 64 行，没人填得动。
实践守则：每个结果节点的直接原因控制在三四个以内，更多因素分层建网。

\subsection{一次完整的贝叶斯诊断计算}

CPT 在手，推理可以双向开弓。
\textbf{预测方向}（由因到果）不费力气：观测到 Lathe\_003 既高温又长时间运行，
查表第一行得故障概率 0.90，超过阈值，触发预防性维护工单。
\textbf{诊断方向}（由果溯因）才见真功夫：维修工接到一台故障设备，
想知道「它处于高温状态的概率多大」，以决定先查冷却系统还是先查别处。
这要用贝叶斯公式反推，包内教学片段给出完整的三步计算，我们逐步带读者手算。

\begin{codebox}
\codeline{\textcolor{cmtgray}{\#\ 第一步，用全概率公式计算故障的边缘概率：}}
\codeline{P(Fault)\ =\ 0.90\(\times\)0.20\(\times\)0.35\ +\ 0.30\(\times\)0.20\(\times\)0.65}
\codeline{\hspace*{9\ttcw}+\ 0.10\(\times\)0.80\(\times\)0.35\ +\ 0.01\(\times\)0.80\(\times\)0.65}
\codeline{\hspace*{9\ttcw}=\ 0.063\ +\ 0.039\ +\ 0.028\ +\ 0.0052}
\codeline{\hspace*{9\ttcw}=\ 0.1352}
\codeblank{}
\codeline{\textcolor{cmtgray}{\#\ 第二步，计算高温条件下的故障概率：}}
\codeline{P(Fault\ |\ HighTemp)\ =\ 0.90\(\times\)0.35\ +\ 0.30\(\times\)0.65\ =\ 0.51}
\codeblank{}
\codeline{\textcolor{cmtgray}{\#\ 第三步，用贝叶斯公式反推：}}
\codeline{P(HighTemp\ |\ Fault)\ =\ P(Fault\ |\ HighTemp)\ \(\times\)\ P(HighTemp)\ /\ P(Fault)}
\codeline{\hspace*{20\ttcw}=\ (0.51\ \(\times\)\ 0.20)\ /\ 0.1352}
\codeline{\hspace*{20\ttcw}\(\approx\)\ 0.754}
\codeblank{}
\codeline{\textcolor{cmtgray}{\#\ 结论：故障设备中约75\%与高温相关\ \(\rightarrow\)\ 维修时优先排查冷却系统}}
\end{codebox}
\color{black}

\assetsource{ch05}{probabilistic-reasoning}

\textbf{第一步：算出「故障」的边缘概率 P(Fault)}。
全概率公式把四种原因组合的贡献加总，每一项 \(=\)
该组合下的故障概率 \(\times\) 该组合自身出现的概率
（两个原因相互独立，组合概率直接相乘）。
逐项验算：高温且长运行贡献 0.90 \(\times\) 0.20 \(\times\) 0.35 \(=\) 0.063；
高温不长运行贡献 0.30 \(\times\) 0.20 \(\times\) 0.65 \(=\) 0.039
（0.65 是「不长运行」的概率，即 1 \(-\) 0.35）；
不高温但长运行贡献 0.10 \(\times\) 0.80 \(\times\) 0.35 \(=\) 0.028；
两者皆无贡献 0.01 \(\times\) 0.80 \(\times\) 0.65 \(=\) 0.0052。
四项相加 0.063 \(+\) 0.039 \(+\) 0.028 \(+\) 0.0052 \(=\) 0.1352——
车间里任意一台设备「正处于故障」的总体概率约 13.5\%。

\textbf{第二步：算 P(Fault \(\mid\) HighTemp)}。
已锁定高温，长运行与否仍未知，按其概率加权两行 CPT：
0.90 \(\times\) 0.35 \(+\) 0.30 \(\times\) 0.65 \(=\) 0.315 \(+\) 0.195 \(=\) 0.51。
高温设备的故障率 51\%，是全体设备 13.5\% 的近四倍——「高温危险」有了定量版本。

\textbf{第三步：贝叶斯反转}。
P(HighTemp \(\mid\) Fault) \(=\) P(Fault \(\mid\) HighTemp) \(\times\) P(HighTemp) \(/\) P(Fault)
\(=\) 0.51 \(\times\) 0.20 \(/\) 0.1352 \(=\) 0.102 \(/\) 0.1352 \(\approx\) 0.754。
解读：在「故障」这个证据到来之前，一台设备高温的先验概率只有 20\%；
证据到来后，后验概率跳到 75.4\%，放大了近 3.8 倍。
维修策略随之而来：\textbf{故障设备四台里约三台与高温相关——先查冷却系统}。
这就是概率推理对调度系统的实际供货：不是「会不会」，而是「先查哪」「值不值得停机」。
建议读者合上书重算一遍三步——所有数字都来自那张 CPT，没有任何黑箱。

算完别忘了审视计算背后的\textbf{建模假设}：
第一步把组合概率直接写成两个先验的乘积，依赖「高温与长运行相互独立」。
现实未必如此——夏季订单高峰时，高温与长运行往往一起出现。
如果台账显示两者相关，正确做法不是去改公式，而是\textbf{改网络结构}：
加一个共同父节点「季节」，或在两个原因之间补一条边，CPT 相应扩行。
这正是贝叶斯网络的纪律性优点：独立性假设全部写在图结构上，
看得见，就改得动；藏在代码或直觉里的假设才最危险。

\subsection{马尔可夫逻辑网络：权重即倾向}

贝叶斯网络要求把变量与依赖结构想清楚，
另一条路线\textbf{马尔可夫逻辑网络}（Markov Logic Network，MLN）则直接给一阶逻辑规则\textbf{标权重}：
权重越大，违反这条规则的世界越「不可能」。
包内教学片段（文件第 4 节）给出三档样例：
权重 2.5 的「高功率设备需要冷却」是强约束，
权重 1.2 的「老旧数控机床会精度漂移」是中等约束，
权重 0.4 的「闲置即可用」只是弱倾向。
权重的语义是指数级的：满足规则的可能世界，概率按 \(e^{w}\) 倍提升——
2.5 对应约 12.2 倍，1.2 对应约 3.3 倍，0.4 只有约 1.5 倍。
权重无穷大时规则退化为不可违反的硬约束——
也就是 OWL 公理「Running \(\sqcap\) Maintenance \(\sqsubseteq\) \(\bot\)」的地位。
MLN 的迷人之处在于\textbf{容忍矛盾}：两条规则打架时不是报告不一致，
而是按权重折中出一个「最不别扭」的世界；
代价是推理变成概率推断（采样或变分近似），结论从「蕴涵」降格为「估计」。
拿样例规则体会折中：一台闲置的老旧数控机床，
权重 0.4 的规则倾向于判它「可用」，权重 1.2 的规则倾向于判它「会精度漂移」。
经典逻辑里若再有「漂移设备不可用」的约束，三者就地爆炸；
MLN 则权衡两边——违反弱规则的代价小，违反中等规则的代价大，
最大概率的世界多半是「可用，但带着漂移风险标记」，
调度器据此把它派给粗加工而非精加工。
\textbf{矛盾不再是故障，而是定价问题}——这就是权重语义的工程价值。

\subsection{工具生态与架构原则}

概率本体的工具生态不如经典推理器成熟，但各路线都有代表作（详见包内教学资产第 5 节）：
\textbf{PR-OWL} 基于多实体贝叶斯网络（MEBN）扩展 OWL；
\textbf{BayesOWL} 把 OWL 类层次编译成贝叶斯网络；
\textbf{Pronto} 是 Pellet 推理器的概率扩展（P-SROIQ 路线）；
\textbf{ProbLog} 与 \textbf{Alchemy} 分别是概率逻辑编程与 MLN 的参考实现。
典型应用场景三类：设备故障诊断与预测性维护（PHM）、
质量风险评估（工艺参数到缺陷概率的映射）、多源传感器数据融合。

工具可以再选，架构原则先立住：

\begin{noteblock}
\textbf{确定性骨架 + 概率叠加。}
类层次、属性定义、不相交等硬约束用 OWL DL 表达，享受可判定推理与一致性保障；
统计性知识（故障率、风险度、置信度）在概率层单独建模，
通过共享的类与个体词汇表对接骨架。
反面教材是把全部知识概率化——硬约束被稀释成「大概率成立」，
一致性检查失去意义，推理复杂度同时失控。
两层分离让每一层都用对自己最划算的数学。
\end{noteblock}

\section{推理性能工程：让推理器扛住生产数据}

样例本体上一切推理都快。真正的考验在生产环境：
几万个类、百万级实例、提交一次推理等四十分钟——
这时拼的不再是逻辑学，而是\textbf{性能工程}。
本节给出四件武器：公理代价意识、模块化抽取、增量推理、物化管道。

\subsection{公理是有价格的}

5.2 节的算法细节在此兑换成一张「价目表」：
不同形态的公理迫使推理器做不同分量的工作，写公理时心里要有这本账。

\begin{center}
\begin{tabularx}{\linewidth}{@{}p{0.28\linewidth}Xp{0.12\linewidth}@{}}
\toprule
\textbf{公理形态} & \textbf{推理器被迫做的事} & \textbf{经验成本} \\
\midrule
原子子类 A \(\sqsubseteq\) B & 层次传播，EL 内核即可消化 & 低 \\
存在限制 \(\exists\)R.C & 造一个证人节点，局部动作 & 低 \\
等价定义 \(\equiv\) & 双向包含：充分性方向要随时反向触发（5.2.3 步骤 1） & 中 \\
全称限制 \(\forall\)R.C & 沿每条 R 边传播标签，与 \(\exists\) 嵌套时逐层级联 & 中高 \\
不相交 / 否定 & 处处埋设冲突检测点，与析取联动放大 & 中高 \\
析取 \(\sqcup\) & 分支加回溯，搜索空间随嵌套指数膨胀 & 高 \\
基数限制 \(\ge\)n / \(\le\)n & 计数、尝试节点合并、合并失败再分支 & 最高 \\
\bottomrule
\end{tabularx}
\end{center}

由此沉淀出与包内教学资产第 5 节一致的改写守则：
\textbf{能用 \(\sqsubseteq\) 不用 \(\equiv\)}——只有真正需要「自动归类」（5.2.3 那种数值触发）的类才配等价定义，
纯分类标签用子类即可；
\textbf{能用 \(\exists\) 不用 \(\forall\)}——「至少有一个」便宜，「全都必须」昂贵，
而且很多 \(\forall\) 写出来其实是误用（想做数据校验的，请回 5.2.6 的提醒）；
\textbf{拆深嵌套}——三层嵌套的复合表达式拆成几个命名中间类，
既给推理器减负，又让分类结果可读；
\textbf{慎用基数}——「恰好 3 个工位」这类约束若只服务于文档目的，降级为注释或 SPARQL 校验。
价目表的另一个用法是\textbf{事后诊断}：推理突然变慢，
先 diff 最近新增的公理里有没有高价条目，命中率惊人。

\subsection{模块化抽取：只带需要的行李}

第二件武器对付本体规模。
调度服务其实只关心设备、任务、状态那几十个类，
却被迫连同整个企业本体（组织、财务、人事）一起推理——冤枉钱花在了无关公理上。
\textbf{模块抽取}（module extraction）解决这件事：
给定一组关心的术语（签名，signature），
从大本体中抽出一个公理子集，并在数学上保证：
\textbf{凡是涉及签名内术语的蕴涵问题，在模块上推理与在全本体上推理，结论完全一致}。
主流算法是基于局部性（locality）的语法抽取，OWL API 内置实现
（\texttt{SyntacticLocalityModuleExtractor}），抽取本身只需秒级。
感受一下量级（数字为示意）：企业本体两万条公理，
调度服务的签名只有 \{Equipment, Task, hasStatus, assignedTo\} 等四十来个术语，
抽出的模块往往只有几百条公理——推理负载直接砍掉两个数量级，
而调度相关的任何蕴涵问题，答案与全本体上完全相同。
工程用法：为每个下游服务按签名抽出专属模块并随服务部署；
本体升级时流水线重新抽取，模块自动跟新——
模块之于本体，恰如视图之于数据库。

\subsection{增量推理：别每次从零开始}

第三件武器对付数据变化频率。
车间数据流是「小步快跑」：每分钟新增几十条状态断言，整库结构纹丝不动。
让推理器每次都全量重算，等于每分钟把三十万类的 SNOMED 重新分类一遍——荒谬，却是默认行为。
\textbf{增量推理}（incremental reasoning）只重算受影响的部分：
新增一条 ABox 断言，只有它能触达的个体与类需要复查，
Pellet 的增量一致性检查、ELK 的增量分类都是此类实现。
要注意\textbf{增删不对称}：增加事实只会增加结论（单调性，第 2 章），增量维护容易；
\textbf{删除}一条事实却可能让一批早先的结论失去依据，
推理器必须追踪「哪些结论依赖哪些事实」才能精确收回，
这正是物化维护算法（如经典的 DRed：先删除受牵连结论，再尝试重推）要解决的难题。
放进车间场景里掂一掂分量：
MES 每分钟写入几十条状态翻转，每条只牵动一台设备的邻域——
增量一致性检查复查这一小片即可，毫秒级完成；
而每条都触发全库重分类的话，推理器永远追不上数据流，队列只会越积越深。
反过来，一旦改动的是 TBox（比如调整了「高功率设备」的功率阈值），
受影响的是\textbf{整类个体}，增量机制省不了多少，老老实实走一次全量批处理更稳妥。
架构启示：\textbf{把高频删改的数据挡在重推理层之外}——
快变状态走 5.3.3 说的查询时轻推理，慢变骨架才交给重型推理器加增量维护；
TBox 变更当作「发版」对待，走批处理加回归验证，而不是混进数据流里偷偷生效。

\subsection{物化回库：与数据管道握手}

第四件武器是把推理从「在线服务」改造成「离线工序」。
判断标准有三条，同时满足就该物化：
\textbf{结论稳定}（TBox 周更而非秒变）、\textbf{查询延迟敏感}（看板和 API 等不起在线推理）、
\textbf{推理深度大}（分类加实现的全家桶，在线算一次太贵）。
做法是把推理挪进数据管道：抽取模块 \(\rightarrow\) 跑推理 \(\rightarrow\)
把推出的三元组写回库 \(\rightarrow\) 下游按普通图查询消费。
两条运维纪律让管道可持续：
其一，\textbf{推理产物与原始事实分图存放}——5.5.2 的具名图在此再立一功，
推出的三元组写进专门的「推理结果图」，可整图丢弃重算，永远不污染人工录入的事实；
其二，\textbf{记录推理器版本与本体版本}作为结果图的元数据，结论可复现、可审计。
这两条纪律合起来保证了管道的\textbf{幂等性}：
任何时刻丢弃结果图、用同版本输入重跑一遍，得到的图逐三元组相同——
出了疑问可以重算对账，迁移存储可以只搬事实图，推理产物随时可再生。
完整的管道编排——调度窗口、失败重试、与 ETL 的衔接——是第 7 章知识图谱管理的正题，
本节的判断标准与分图纪律是那套管道的推理侧约定。

\section{本章小结}

本章从服务、算法、机制三个层面打开了推理器的黑箱。
\textbf{服务层}：分类、一致性检查、可满足性、实例检查与实现、查询回答五种标准服务，
都是「KB 是否蕴涵」的不同问法，前向与后向是付推理成本的两种时机。
\textbf{算法层}：Tableau 用「试着造世界」回答可满足性，
四条展开规则配合分支回溯，阻塞保证停机；
consequence-based 的 ELK 用砍表达力换多项式速度，RETE 用缓存换增量匹配，
物化与查询时推理是前向与后向在架构上的投影。
\textbf{机制层}：SWRL 在 DL-Safe 约束下把多变量业务规则写进本体，但写不了否定；
具名图给三元组打时间戳，Allen 十三式是时态推理的坐标系；
概率标注与贝叶斯计算让「大概率」可推理，MLN 把硬约束放松成权重；
性能工程的四件武器——公理价目表、模块抽取、增量推理、物化管道——
让这一切扛得住生产数据。
第 2 章的逻辑地基、第 4 章的 OWL 语法、本章的推理机制至此会师；
第 7 章将把 Dataset、SHACL 与版本治理接起来，第 9 章则展示规则资产存在、
受限 profile 可执行与一般 SWRL 仍阻断怎样在同一 package 中同时诚实表达。

给动手派读者一条练习路线：
先对照 5.2.3 与包内教学资产把 Lathe\_003 的推演\textbf{手算}一遍（含归谬视角）；
再查看 ch05 manifest，运行 \texttt{OE-V1-CH05-SCN-FORWARD-CHAIN-001}，
核对两条 exact conclusion；随后尝试请求未声明的 SWRL/DL/时态/概率操作，确认 runner
fail closed。最后比较手算理论与受限执行 profile：知道系统没做什么，与知道它做了什么同样重要。

\subsection*{思考题}

\begin{enumerate}
\item 沿用 5.2.3 的 TBox，ABox 改为「设备(Mill\_007)、功率(Mill\_007, 9.8)」。
  用 Tableau 步骤判断：推理器会得出「Mill\_007 需要冷却」吗？
  进一步，推理器能断言「Mill\_007 \textbf{不是}高功率设备」吗？
  提示：第二问与开放世界假设有关——9.8 是它唯一的功率值吗？知识库说了吗？
\item 对循环公理「设备 \(\sqsubseteq\) \(\exists\)属于.产线」「产线 \(\sqsubseteq\) \(\exists\)包含.设备」，
  从个体 Lathe\_003 出发手画完成图，逐节点写出标签，标注阻塞在第几个节点触发、为什么。
\item 用 5.4 的词汇表写一条 SWRL 规则：同一设备上相邻两道工序，
  若前道结束时间与后道开始时间间隔小于 600 秒，则给后道工序贴「换刀检查」标签。
  写出所需的全部属性原子与内建原子，并说明变量绑定顺序为什么不能颠倒。
\item 对照 Allen 十三式，验证 5.4.4 时间冲突规则的双不等式恰好命中「存在公共内部」的九种关系；
  若把两个严格不等号都改成「大于等于／小于等于」，会额外命中哪两种关系？
  这对「首尾相接的流水工序」意味着什么？
\item 夏季车间高温更频繁：把 5.6 的先验 P(HighTemp) 从 0.20 改为 0.50，
  重算 P(Fault) 与 P(HighTemp \(\mid\) Fault)。
  对比冬夏两组数字，说明「先验变化会不会改变 CPT」以及「维修排查策略该不该换季调整」。
\item 业务需求「列出所有未分配任务的设备」。分别用 5.4.5 的三条绕法给出实现思路，
  比较三者的数据维护成本、查询复杂度与语义风险，说明你在什么团队条件下选哪条。
\item 某调度系统每天新增约 5 万条状态三元组，查询峰值 200 QPS 且要求百毫秒级延迟。
  用 5.3.3 的对照表与 5.7 的四件武器，设计推理架构：
  哪些知识物化、哪些查询时算、增量还是全量、推理结果存进哪个图？
\item 回看第 2 章「高端设备」的定义：
  \begin{center}
  HighEndEquipment \(\equiv\) Equipment \(\sqcap\) (price \(\ge\) 1000000)
  \(\sqcap\) (\(\forall\)hasFeature.AdvancedFeature)
  \end{center}
  它在 5.7.1 价目表上命中哪几条高价条目？
  给出两种降低推理成本的改写方案，并说明各自损失了什么语义。
\end{enumerate}

\subsection*{本章包资产与场景指引}

本章全部片段来自 ch05 包内下列资产；建议对照正文阅读，
注释里还有不少正文未及展开的细节。

\begin{center}
\begin{tabularx}{\linewidth}{@{}>{\ttfamily\footnotesize}p{0.40\linewidth}X@{}}
\toprule
\normalfont\textbf{Semantica asset/scenario} & \textbf{内容与阅读建议} \\
\midrule
description-logic-examples & Tableau 直觉、手算推演与一致性反例；教学 case，不是 DL runner \\
swrl-rules & 五类调度规则与语法；来源保留，一般 SWRL/built-ins 不执行 \\
temporal-reasoning & TriG 快照与时态规则；教学材料，无执行 oracle \\
probabilistic-reasoning & 概率公理、CPT 与贝叶斯计算；教学材料，无执行 oracle \\
OE-V1-CH05-SCN-FORWARD-CHAIN-001 & 受限正向链 exact oracle；不代表上述四类高级 profile 已支持 \\
\bottomrule
\end{tabularx}
\end{center}
